\scanpage{62}

c) Pour tt $\omega\in\operatorname{Or}(\Pi)$, $\exists\unsure{1|!}u_\omega\in G^+$ tel que $u_\omega\mid R_\omega^+=u\mid R_\omega^+$, et on a \deleted{\unsure{$u_\omega$}}

\[u_{-\omega}=u_\omega^{-1}\]

\deleted{Donc} De plus, $u_\omega$ est un générateur de $G^+$, et si $n\geq 3$, $\omega\mapsto u_\omega$ est une injection.

\[\operatorname{Or}(\Pi)\unsure{\leftrightarrow|\hookrightarrow}\operatorname{Génér.}(G^+)\]

\margin{\deleted{\unsure{$\nearrow$}}\ $\nearrow$\ $\nearrow$}
% slanted

d) Le foncteur

\[\Pi\longmapsto (Z=G^+,\inserted{\Omega=}\{u_\omega,u_{-\omega}\}\in G^+,S)\]

des $n$-polygones \inserted{$(\unsure{3|{-3}}\leq{n}\leq+\infty)$} vers les triples d'un groupe cyclique d'ordre $n$ $G^+$, d'un ens de deux gén. de $G^+$ inv. l'un de l'autre, et d'un $G^+$ torseur $S$, est une équivalence de \unsure{ctégories|catégories}.

\uline{\unsure{Exer|Exo|Exer.}} Exprimer $A$ \inserted{(en \unsure{tt|tant} que $Z$-torseur)} en termes des données précédentes, ainsi que $R$ en tt que partie de $S\times R$ \dots

% archive sheet number: 61
